% =========================================================
% ESTILO GERAL — TEMPLATE UFRN 2026
% Compilador: pdfLaTeX
% =========================================================

% PACOTES DE FONTES E IDIOMA
\usepackage{fix-cm}
\usepackage{ragged2e}
\usepackage[utf8]{inputenc}
\usepackage[T1]{fontenc}
\usepackage[brazilian]{babel}

% FONTES
% Definicoes de fontes guardadas em modelo/fontes/definicoes/.
\ProvidesFile{ot1lbr.fd}[2026/10/12 Liberation Sans OT1 fallback]
\DeclareFontFamily{OT1}{lbr}{}
\DeclareFontShape{OT1}{lbr}{m}{n}{<->ssub*cmss/m/n}{}
\DeclareFontShape{OT1}{lbr}{bx}{n}{<->ssub*cmss/bx/n}{}
\DeclareFontShape{OT1}{lbr}{b}{n}{<->ssub*cmss/bx/n}{}
\DeclareFontShape{OT1}{lbr}{m}{it}{<->ssub*cmss/m/sl}{}
\DeclareFontShape{OT1}{lbr}{m}{sl}{<->ssub*cmss/m/sl}{}
\endinput

\ProvidesFile{t1lbr.fd}[2026/10/12 Liberation Sans font family]
\DeclareFontFamily{T1}{lbr}{}
\DeclareFontShape{T1}{lbr}{m}{n}{<-> modelo/fontes/lbrr}{}
\DeclareFontShape{T1}{lbr}{b}{n}{<-> modelo/fontes/lbrb}{}
\DeclareFontShape{T1}{lbr}{bx}{n}{<-> modelo/fontes/lbrb}{}
\DeclareFontShape{T1}{lbr}{m}{it}{<-> modelo/fontes/lbri}{}
\DeclareFontShape{T1}{lbr}{m}{sl}{<-> modelo/fontes/lbri}{}
\DeclareFontShape{T1}{lbr}{b}{it}{<-> modelo/fontes/lbrbi}{}
\DeclareFontShape{T1}{lbr}{bx}{it}{<-> modelo/fontes/lbrbi}{}
\endinput

\ProvidesFile{ts1lbr.fd}[2026/11/05 Text symbols (TS1) for Liberation Sans]
\DeclareFontFamily{TS1}{lbr}{}
\DeclareFontShape{TS1}{lbr}{m}{n}{<->ssub*phv/m/n}{}
\DeclareFontShape{TS1}{lbr}{b}{n}{<->ssub*phv/b/n}{}
\DeclareFontShape{TS1}{lbr}{bx}{n}{<->ssub*phv/b/n}{}
\DeclareFontShape{TS1}{lbr}{m}{it}{<->ssub*phv/m/sl}{}
\DeclareFontShape{TS1}{lbr}{m}{sl}{<->ssub*phv/m/sl}{}
\DeclareFontShape{TS1}{lbr}{b}{it}{<->ssub*phv/b/sl}{}
\DeclareFontShape{TS1}{lbr}{bx}{it}{<->ssub*phv/b/sl}{}
\endinput


% Liberation Sans incluida no projeto.
\pdfmapfile{+modelo/fontes/liberation.map}
\renewcommand{\sfdefault}{lbr}
\renewcommand{\rmdefault}{lbr}
\renewcommand{\familydefault}{\sfdefault}

% MARGENS
\usepackage{geometry}
\geometry{a4paper,left=3cm,top=3cm,right=2cm,bottom=2cm,headheight=15pt,headsep=0.7cm}

% ESPACAMENTO E PARAGRAFOS
\usepackage{setspace}
\usepackage{indentfirst}
\onehalfspacing
\setlength{\parindent}{1.5cm}
\setlength{\parskip}{0pt}

% PACOTES GERAIS
\usepackage{graphicx}
% Imagens em 5-figuras/.
\graphicspath{{5-figuras/}}
\usepackage{array,booktabs,tabularx,longtable}
\usepackage{float}
\usepackage{microtype}
\usepackage{caption}
\usepackage{etoolbox}
\usepackage{enumitem}
\usepackage{amsmath,amssymb}
% tikz e pgfplots desativados: nenhum arquivo do projeto os utiliza
% (todas as figuras sao PNG/PDF externos) e custam tempo de compilacao.
% \usepackage{tikz}
% \usepackage{pgfplots}
% \pgfplotsset{compat=1.18}
\usepackage{pdfpages}
\usepackage{fancyhdr}
\usepackage{titlesec}
\usepackage{tocloft}
\usepackage[hyphens]{url}

% =========================================================
% REFERENCIAS AUTOMATICAS — ABNTEX2CITE / BIBTEX
% O arquivo .bib fica na pasta 4-referencias.
% =========================================================
\usepackage{csquotes}
% PDF/A-2b (carrega o hyperref). Titulo e autor vem de 1-dados/dados.tex,
% que precisa ser lido antes daqui (\CarregarDados em Principal.tex).
% O pdfx le <nome da compilacao>.xmpdata (Principal.xmpdata localmente;
% "output.xmpdata" no Overleaf), gerado abaixo quando nao existe.
\makeatletter
\ifdefined\titulo\else
  \PackageError{Modelo PPGENAP}{Dados nao carregados antes de modelo/estilo.tex}%
    {Restaure Principal.tex a partir do ZIP do modelo: os metadados PDF/A ficariam sem titulo e autor.}%
\fi
\IfFileExists{\jobname.xmpdata}{}{%
  \newwrite\ppg@xmp
  \immediate\openout\ppg@xmp=\jobname.xmpdata\relax
  \immediate\write\ppg@xmp{\string\Title{\string\titulo}}%
  \immediate\write\ppg@xmp{\string\Author{\string\autor}}%
  \immediate\write\ppg@xmp{\string\Language{pt-BR}}%
  \immediate\closeout\ppg@xmp}
\makeatother
\usepackage[a-2b]{pdfx}
\IfFileExists{\jobname.xmpdata}{}{%
  \PackageError{Modelo PPGENAP}{Metadados PDF/A (\jobname.xmpdata) nao encontrados}%
    {O PDF ficou sem titulo e autor nos metadados. Compile novamente; o arquivo e gerado automaticamente.}}
\hypersetup{hidelinks}
% Opcoes ABNT em modelo/abnt-opcoes.bib.
\usepackage[alf,abnt-emphasize=bf,abnt-ufrn]{abntex2cite}
\providecommand{\textcite}{\citeonline}
\providecommand{\citet}{\citeonline}
% "apud" em italico, como o "et al." (ver abnt-etal-text em abnt-opcoes.bib).
\renewcommand{\apudname}{\textit{apud}}
% Citacao de citacao com pagina da obra original e/ou da consultada
% (deixe {} vazio para omitir uma delas):
%   \apudpag{original}{pag. original}{consultada}{pag. consultada}
%   \apudpagonline{original}{pag. original}{consultada}{pag. consultada}
% Trecho comum: "ano[, pag.] apud Autor, ano[, pag.])".
\makeatletter
\newcommand{\ppg@apudpag@resto}[4]{\citeyear{#1}%
\ifthenelse{\equal{#2}{}}{}{, #2} \apudname\ %
\citeauthor{#3}, \citeyear{#3}%
\ifthenelse{\equal{#4}{}}{}{, #4})}
\newcommand{\apudpag}[4]{(\citeauthor{#1}, \ppg@apudpag@resto{#1}{#2}{#3}{#4}}
\newcommand{\apudpagonline}[4]{\citeauthoronline{#1} (\ppg@apudpag@resto{#1}{#2}{#3}{#4}}
\makeatother
\bibliographystyle{modelo/abntex2-alf-ufrn}
\def\UrlLeft{}\def\UrlRight{}
% Autor institucional no .bib: author = {{\autorinstitucional{Nome}}}
\DeclareRobustCommand{\autorinstitucional}[1]{#1}
\addto\captionsbrazilian{%
  \renewcommand{\bibname}{REFERÊNCIAS}%
}

% Bibliografia: espacamento simples, alinhada a esquerda,
% sem recuo inicial, com separacao entre referencias.
\AfterTheBibliography{%
  \DeclareRobustCommand{\autorinstitucional}[1]{\MakeUppercase{##1}}%
  \singlespacing
  \raggedright
  \setlength{\itemsep}{\baselineskip}%
  \phantomsection
  \addcontentsline{toc}{chapter}{REFERÊNCIAS}%
}


% TITULOS DOS CAPITULOS E SECOES
% Modelo BCZM/UFRN: um espaco simples entre numero e titulo; uma linha
% em branco (entrelinha 1,5) depois do titulo primario e antes e depois
% dos demais titulos.
\newcommand{\ppgsepnumero}{0.3em}
\titleformat{\chapter}[hang]{\normalfont\bfseries\normalsize}{\thechapter}{\ppgsepnumero}{}
\titlespacing*{\chapter}{0pt}{-0.78cm}{\baselineskip}
\titleformat{name=\chapter,numberless}[hang]{\normalfont\bfseries\normalsize\filcenter}{}{0pt}{}
\titleformat{\section}{\normalfont\normalsize\bfseries}{\thesection}{\ppgsepnumero}{}
\titlespacing*{\section}{0pt}{\baselineskip}{\baselineskip}
\titleformat{\subsection}{\normalfont\normalsize}{\thesubsection}{\ppgsepnumero}{}
\titlespacing*{\subsection}{0pt}{\baselineskip}{\baselineskip}
\titleformat{\subsubsection}{\normalfont\normalsize\itshape}{\thesubsubsection}{\ppgsepnumero}{}
\titlespacing*{\subsubsection}{0pt}{\baselineskip}{\baselineskip}
\titleformat{\paragraph}{\normalfont\normalsize}{\theparagraph}{\ppgsepnumero}{}
\titlespacing*{\paragraph}{0pt}{\baselineskip}{\baselineskip}

% TITULOS PRE-TEXTUAIS
\newcommand{\pretitle}[1]{\chapter*{\MakeUppercase{#1}}}

% LEGENDAS E TEXTOS DE EXEMPLO
% \fonte: obrigatoria, logo abaixo da ilustracao ou tabela.
% \legenda: opcional, abaixo da fonte (explicacao de cores, simbolos etc.).
\newcommand{\ppgtextoilustracao}[2]{{\fontsize{11}{13.2}\selectfont\setstretch{1}\raggedright\noindent #1: #2\par}}
\newcommand{\fonte}[1]{\par\vspace{4pt}\ppgtextoilustracao{Fonte}{#1}}
\newcommand{\legenda}[1]{\par\ppgtextoilustracao{Legenda}{#1}}
\newcommand{\textopadrao}{Texto texto texto texto texto texto texto texto texto texto texto texto texto texto texto texto texto texto texto texto texto texto texto texto.}

% CITACOES LONGAS
\newenvironment{citacaolonga}{%
  \par\vspace{\baselineskip}%
  \begin{list}{}{\setlength{\leftmargin}{4cm}\setlength{\rightmargin}{0cm}%
    \setlength{\topsep}{0pt}\setlength{\partopsep}{0pt}%
    \setlength{\parsep}{0pt}\setlength{\itemsep}{0pt}}%
  \item\fontsize{11}{13.2}\selectfont\singlespacing\noindent
}{\end{list}\vspace{\baselineskip}}

% FIGURAS, TABELAS, QUADROS E GRAFICOS
% Modelo BCZM/UFRN: titulo centralizado acima; fonte e legenda abaixo,
% alinhadas a margem esquerda da ilustracao.
\captionsetup{font={normalsize,singlespacing},labelfont=normalfont,justification=centering,singlelinecheck=false,labelsep=endash,position=top,skip=6pt}
% Uso: \begin{ilustracao}{0.8\linewidth} ... \end{ilustracao}
\newenvironment{ilustracao}[1]{%
  \par\centering\begin{minipage}{#1}\setlength{\parindent}{0pt}%
}{%
  \end{minipage}\par
}
% Cabecalho de coluna em negrito e centralizado (quadros e tabelas).
\newcommand{\titulocoluna}[1]{\centering\arraybackslash\textbf{#1}}
\captionsetup[table]{position=top}
\renewcommand{\thefigure}{\arabic{figure}}
\renewcommand{\thetable}{\arabic{table}}
\makeatletter
\@removefromreset{figure}{chapter}
\@removefromreset{table}{chapter}
\makeatother
\usepackage{newfloat}
\DeclareFloatingEnvironment[fileext=loq,name=Quadro,placement=htbp,within=none]{quadro}
\DeclareFloatingEnvironment[fileext=lgf,name=Gráfico,placement=htbp,within=none]{grafico}
\captionsetup[quadro]{position=top}
\captionsetup[grafico]{position=top}
\renewcommand{\thequadro}{\arabic{quadro}}
\renewcommand{\thegrafico}{\arabic{grafico}}

% LISTAS E SUMARIO
\renewcommand{\listfigurename}{LISTA DE FIGURAS}
\renewcommand{\listtablename}{LISTA DE TABELAS}
\renewcommand{\contentsname}{SUMÁRIO}
\addto\captionsbrazilian{%
 \renewcommand{\listfigurename}{LISTA DE FIGURAS}%
 \renewcommand{\listtablename}{LISTA DE TABELAS}%
 \renewcommand{\contentsname}{SUMÁRIO}%
}
\renewcommand{\cftloftitlefont}{\hfill\normalsize\bfseries}
\renewcommand{\cftafterloftitle}{\hfill}
\renewcommand{\cftlottitlefont}{\hfill\normalsize\bfseries}
\renewcommand{\cftafterlottitle}{\hfill}
% A centralizacao pelo \hfill deslocava o titulo para a direita.
% Centraliza na pagina, compensando margens de 3cm e 2cm.
\renewcommand{\cfttoctitlefont}{\normalfont\normalsize\bfseries}
\renewcommand{\cftaftertoctitle}{}
\makeatletter
\renewcommand{\@cftmaketoctitle}{%
  \addpenalty\@secpenalty
  \vspace*{\cftbeforetoctitleskip}%
  \@cftpagestyle
  {\interlinepenalty\@M
    \noindent\hspace*{-0.5cm}\makebox[\linewidth][c]{\cfttoctitlefont\contentsname}%
    \cftmarktoc
    \par\nobreak
    \vskip\cftaftertoctitleskip
    \@afterheading}%
}
\makeatother
\setlength{\cftbeforetoctitleskip}{-0.78cm}
\setlength{\cftaftertoctitleskip}{1.25\baselineskip}
\setlength{\cftbeforeloftitleskip}{-0.78cm}
\setlength{\cftbeforelottitleskip}{-0.78cm}

% FORMATACAO AUTOMATICA DO SUMARIO
% Aplicada automaticamente a capitulos e secoes futuros.
% O agrupamento limita o espacamento simples somente ao sumario.
\let\sumarioOriginal\tableofcontents
\renewcommand{\tableofcontents}{%
  \begingroup
  \singlespacing%
  \setlength{\parindent}{0pt}%
  \sumarioOriginal
  \endgroup
}
\setlength{\cftbeforechapskip}{2.7pt}
\setlength{\cftbeforesecskip}{2.7pt}
\setlength{\cftbeforesubsecskip}{2.7pt}
\setlength{\cftbeforesubsubsecskip}{2.7pt}
\setlength{\cftparskip}{0pt}

% MODELO WORD UFRN — ALINHAMENTO DAS ENTRADAS
% Os numeros de todos os niveis comecam na mesma margem esquerda.
% A largura reservada para cada numero determina onde o titulo comeca.
% Nao acrescente outro bloco de recuos no final do arquivo.
\setlength{\cftchapindent}{3pt}
\setlength{\cftchapnumwidth}{2cm}
\setlength{\cftsecindent}{3pt}
\setlength{\cftsecnumwidth}{2cm}
\setlength{\cftsubsecindent}{3pt}
\setlength{\cftsubsecnumwidth}{2cm}
\setlength{\cftsubsubsecindent}{3pt}
\setlength{\cftsubsubsecnumwidth}{2cm}

% Destaque dos niveis
\renewcommand{\cftchapfont}{\bfseries}
\renewcommand{\cftchappagefont}{\normalfont}
\renewcommand{\cftsecfont}{\bfseries}
\renewcommand{\cftsecpagefont}{\normalfont}
\renewcommand{\cftsubsecfont}{\normalfont}
\renewcommand{\cftsubsecpagefont}{\normalfont}
\renewcommand{\cftsubsubsecfont}{\itshape}
\renewcommand{\cftsubsubsecpagefont}{\normalfont}

% Pontilhados em todos os niveis (incluindo capitulos sem numero)
\renewcommand{\cftdotsep}{0.5}
\renewcommand{\cftchapleader}{\cftdotfill{\cftdotsep}}
\renewcommand{\cftsecleader}{\cftdotfill{\cftdotsep}}
\renewcommand{\cftsubsecleader}{\cftdotfill{\cftdotsep}}
\renewcommand{\cftsubsubsecleader}{\cftdotfill{\cftdotsep}}


% Entradas sem numero (REFERENCIAS, APENDICES e ANEXOS)
% ficam na mesma coluna do titulo dos capitulos numerados.
\makeatletter
\ExplSyntaxOn
\cs_new_eq:NN \ufrn_orig_lchapter:nn \l@chapter
\cs_set_protected:Npn \l@chapter #1#2
  {
    \exp_args:Nee \str_if_in:nnTF
      {\tl_to_str:n {#1}} {\tl_to_str:n {\numberline}}
      {\ufrn_orig_lchapter:nn {#1} {#2}}
      {\ufrn_orig_lchapter:nn {\numberline{}#1} {#2}}
  }
\ExplSyntaxOff
\makeatother

% IDENTIFICACAO DAS LISTAS
\setlength{\cftfigindent}{0pt}
\setlength{\cfttabindent}{0pt}
\renewcommand{\cftfigpresnum}{Figura }
\renewcommand{\cftfigaftersnum}{ – }
\setlength{\cftfignumwidth}{5.2em}
\renewcommand{\cfttabpresnum}{Tabela }
\renewcommand{\cfttabaftersnum}{ – }
\setlength{\cfttabnumwidth}{5.5em}

% LISTAS DE QUADROS E GRAFICOS
\providecommand{\listofquadros}{}
\providecommand{\listofgraficos}{}
\newcommand{\titulolista}[1]{\addcontentsline{toc}{chapter}{#1}}
\makeatletter
\renewcommand{\listofgraficos}{%
  \clearpage\chapter*{LISTA DE GRÁFICOS}%
  \@starttoc{lgf}}
\renewcommand{\listofquadros}{%
  \clearpage\chapter*{LISTA DE QUADROS}%
  \@starttoc{loq}}
\providecommand*{\l@grafico}[2]{}
\renewcommand*{\l@grafico}[2]{%
  \begingroup
  \let\oldnumberline\numberline
  \renewcommand{\numberline}[1]{\oldnumberline{Gráfico~##1\,–}}%
  \@dottedtocline{1}{0em}{6.1em}{#1}{#2}\endgroup}
\providecommand*{\l@quadro}[2]{}
\renewcommand*{\l@quadro}[2]{%
  \begingroup
  \let\oldnumberline\numberline
  \renewcommand{\numberline}[1]{\oldnumberline{Quadro~##1\,–}}%
  \@dottedtocline{1}{0em}{6.1em}{#1}{#2}\endgroup}
\makeatother

% NUMERACAO DAS PAGINAS
\newcommand{\numeropagina}{{\fontsize{11}{13.2}\selectfont\thepage}}
\newcommand{\paginassemnumero}{%
  \pagestyle{empty}%
  \fancypagestyle{plain}{\fancyhf{}\renewcommand{\headrulewidth}{0pt}}%
}
\newcommand{\paginascomnumero}{%
  \pagestyle{fancy}%
  \fancypagestyle{plain}{\fancyhf{}\fancyhead[R]{\numeropagina}\renewcommand{\headrulewidth}{0pt}}%
}
\pagestyle{fancy}
\fancyhf{}
\fancyhead[R]{\numeropagina}
\renewcommand{\headrulewidth}{0pt}
\fancypagestyle{plain}{\fancyhf{}\fancyhead[R]{\numeropagina}\renewcommand{\headrulewidth}{0pt}}
\setlength{\footskip}{24pt}
\setlength{\topskip}{12pt}
\setcounter{secnumdepth}{4}
\setcounter{tocdepth}{4}

% NOTAS DE RODAPE
\makeatletter
\renewcommand\footnotesize{\@setfontsize\footnotesize{11pt}{13.2pt}}
\makeatother
\usepackage{footmisc}
\renewcommand{\footnotelayout}{\setstretch{1}}

% REFERENCIAS
\newenvironment{referenciasufrn}{%
  \begingroup\singlespacing\raggedright\setlength{\parindent}{0pt}%
  \setlength{\parskip}{\baselineskip}%
}{\par\endgroup}

% NATUREZA DO TRABALHO
\newcommand{\naturezadotrabalho}{%
  \setstretch{1}%
  \tipotrabalho\ apresentada ao curso de
  \nivelcurso\ em \nomecurso,
  da Universidade Federal do Rio Grande do Norte,
  como requisito parcial à obtenção do título de
  \grau\ em \areagrau.
}

% APENDICES E ANEXOS
% Uso: \apendice{TÍTULO DO APÊNDICE} e \anexo{TÍTULO DO ANEXO}.
% A letra (A, B, C...) e a entrada no sumario sao automaticas.
\newcounter{apendice}
\renewcommand{\theapendice}{\Alph{apendice}}
\newcommand{\apendice}[1]{%
  \clearpage\refstepcounter{apendice}%
  \chapter*{APÊNDICE \theapendice\ --- #1}%
  \phantomsection
  \addcontentsline{toc}{chapter}{APÊNDICE \theapendice\ --- #1}%
}
\newcounter{anexo}
\renewcommand{\theanexo}{\Alph{anexo}}
\newcommand{\anexo}[1]{%
  \clearpage\refstepcounter{anexo}%
  \chapter*{ANEXO \theanexo\ --- #1}%
  \phantomsection
  \addcontentsline{toc}{chapter}{ANEXO \theanexo\ --- #1}%
}
